%% kapa: exam=1 round=24 year=2013 period=2 form=A subject=appraisal_law subject_ko=부동산관계법규 seq=33 number=73
%% source: Q-Net 감정평가사 기출문제 (cst00309) artlSeq=5073802 file=extracted/1008/5073802/2013년 제24회 감정평가사 1차시험 기출문제/2교시(A형).hwp
%% conversion: kordoc 4.24.51 -> tools/gen1.py (auto); review=proofread-sonnet; flags=none
\begin{question}{73}
측량ㆍ수로조사 및 지적에 관한 법령상 지상 경계를 새로 결정하려는 경우의 기준으로 옳은 것은? (단, 지상 경계의 구획을 형성하는 구조물 등의 소유자가 다르지 않은 경우를 전제로 한다.)
\begin{choices}
  \item 연접되는 토지 간에 높낮이 차이가 없는 경우: 그 구조물 등의 중앙
  \item 연접되는 토지 간에 높낮이 차이가 있는 경우: 그 구조물 등의 상단부
  \item 도로ㆍ구거 등의 토지에 절토(切土)된 부분이 있는 경우: 그 경사면의 하단부
  \item 토지가 해면 또는 수면에 접하는 경우: 최소만조위 또는 최소만수위가 되는 선
  \item 공유수면매립지의 토지 중 제방 등을 토지에 편입하여 등록하는 경우: 안쪽 어깨부분
\end{choices}
\end{question}
\answer{1}
